\section{Seguimiento y registro de hiperparámetros}

\subsection{Estructura del logbook}

El curso recomienda mantener un logbook para cada experimento, en el que se registren los siguientes campos:

\begin{itemize}
    \item \textbf{Learner IDs}: la LR de Actor, Critic y Target network;
    \item \textbf{Replay ratio}: las transitions consumidas por cada update;
    \item \textbf{Entropy schedule}: el valor inicial más la velocidad de decay;
    \item \textbf{Benchmark metrics}: KL, critic loss, reward mean.
\end{itemize}

\begin{importantbox}{Por qué log mutation}
Al comparar distintos seeds, con frecuencia solo hay una ligera desviación en los hiperparámetros. Llevar un logbook completo te permite reproducir rápidamente un good run, y también detectar rápidamente una catastrophic divergence.
\end{importantbox}

\subsection{Tablas automatizadas}

\begin{table}[H]
\centering
\begin{tabular}{p{0.25\textwidth}p{0.25\textwidth}p{0.25\textwidth}}
\toprule
Campo & Herramienta recomendada & Objetivo \\
\midrule
KL curve & TensorBoard / WandB & Monitorear si se rebasa el trust region \\
Entropy & WandB custom scalar & Dar seguimiento al exploration collapse \\
Critic loss & CSV + daily summary & Detectar señales tempranas de critic drift \\
Replay ratio & elastic log & Asegurar que no se olvide la sample efficiency \\
\bottomrule
\end{tabular}
\caption{Prácticas concretas de Hyperparameter logging}
\end{table}

\subsection{Resumen de la sección}

El registro de hiperparámetros eleva el experimento de «ejecutar algunos seeds» a «un proceso de ingeniería reproducible», y lo alinea con la colaboración en equipo y la revisión de performance regression.

